\section{Proofs of spectral preservation and statistical results}
\label{app:statistical-proofs}

\subsection{Exact preservation and the population contrast space}

\begin{proof}[Proof of Proposition~\ref{prop:exact-preservation}]
Since $AU=U\Lambda$, multiplication by $U^\top$ gives
$A\Pi=AUU^\top=U\Lambda U^\top=\widetilde A$.  Therefore
$\widetilde A_iw-A_iw=A_i(\Pi-\mathrm{Id})w$.  If $w\in\operatorname{col}(U)$,
then $\Pi w=w$ and the difference vanishes.

Let $N=Z^\top Z=\operatorname{diag}(n_1,\ldots,n_K)$, which is
nonsingular since all communities are nonempty.  The columns of
$ZN^{-1/2}$ are orthonormal, and
\[
 ZPZ^\top=(ZN^{-1/2})(N^{1/2}PN^{1/2})(ZN^{-1/2})^\top.
\]
Nonsingularity of $P$ makes the central matrix nonsingular.  Hence
$ZPZ^\top$ has rank $K$ and column space $\operatorname{col}(Z)$.
For each contrast, define
$\xi_c^{ab}=\log\{P_{ac}(1-P_{bc})/[P_{bc}(1-P_{ac})]\}$.
Then $w^{ab}=Z\xi^{ab}$, which proves membership in that space.
Finally, subtracting $\operatorname{diag}(P_{z_i z_i})$ from $ZPZ^\top$
gives the actual no-self-loop mean, so no rank assertion about that
corrected matrix is needed.
\end{proof}

Full rank cannot be replaced by distinct block rows without a further
contrast-span condition.  For example, let $P=(\log n/n)ss^\top$ for a positive,
nonconstant vector $s$.  Its signal space is spanned by $Zs$, whereas
the sparse log-probability contrast between communities $a,b$ is
$\log(s_a/s_b)\mathbf1$.  When $s_a\ne s_b$, this constant vector does
not belong to $\operatorname{span}(Zs)$.  A rank-one signal projection
does not preserve that contrast by the preceding identity.  This example
only concerns projection onto the population rank-one space; it is not
an impossibility claim for every statistic computed from one empirical
eigenvector.

\subsection{The sparse likelihood correction}

\begin{proof}[Proof of Lemma~\ref{lem:sparse-expansion}]
Write $p_{aj}=P_{a,z_j}$ and $p_{bj}=P_{b,z_j}$ for
$j\ne i$.  The exact binary Bernoulli log-likelihood ratio is
\[
 L_i^{ab}=\sum_{j\ne i}A_{ij}\log\frac{p_{aj}}{p_{bj}}
 +\sum_{j\ne i}A_{ij}\log\frac{1-p_{bj}}{1-p_{aj}}
 +\sum_{j\ne i}\log\frac{1-p_{aj}}{1-p_{bj}}.
\]
For $0\le x\le1/2$, the power series for $\log(1-x)$ implies
$|\log(1-x)+x|\le Cx^2$.  Since
$p_{aj},p_{bj}\le b_+\log n/n$, it also implies
$|\log((1-p_{bj})/(1-p_{aj}))|\le C\log n/n$ and
\[
 \left|\log\frac{1-p_{aj}}{1-p_{bj}}+(p_{aj}-p_{bj})\right|
 \le C(\log n)^2/n^2.
\]
The first term is unchanged in the sparse score, and the final term is replaced by its first-order nonedge offset.
Summing the two displayed remainders proves \eqref{eq:sparse-expansion-error}.

Under any true community, $\deg(i)$ is a sum of independent Bernoulli
variables with mean at most $b_+\log n$.  For $D>b_+$, exponential Markov
inequality with $t=\log(D/b_+)$ and the bound
$\log\mathbb E e^{t\deg(i)}\le b_+\log n(e^t-1)$ give
\[
 \Pr\{\deg(i)>D\log n\}
 \le\exp\{-\log n[D\log(D/b_+)-D+b_+]\}.
\]
The bracket tends to infinity with $D$, so $D_H$ can be chosen for any
$H$.  On the complementary event the correction is at most
$C(D_H+1)(\log n)^2/n=o(1)$.
\end{proof}

\subsection{Chernoff information and the oracle exponent}

\begin{proof}[Proof of Lemma~\ref{lem:oracle-exponent}]
Let $f_a,f_b$ denote the product Bernoulli mass functions of the
incident-edge row in the two hypotheses.  Define
\begin{align*}
 D_\alpha(p_{a,-i}\Vert p_{b,-i})
 &=-\log\sum_x f_a(x)^{1-\alpha}f_b(x)^\alpha\\
 &=-\sum_{j\ne i}\log\big\{p_{aj}^{1-\alpha}p_{bj}^\alpha
                    +(1-p_{aj})^{1-\alpha}(1-p_{bj})^\alpha\big\},\\
 D^*_{i,ab}&=\max_{0\le \alpha\le1}D_\alpha(p_{a,-i}\Vert p_{b,-i}).
\end{align*}
Uniformly for $\alpha\in[0,1]$ and
$u,v\in[b_-\log n/n,b_+\log n/n]$,
\[
 u^{1-\alpha}v^\alpha+(1-u)^{1-\alpha}(1-v)^\alpha
 =1-\{(1-\alpha)u+\alpha v-u^{1-\alpha}v^\alpha\}
   +O((\log n)^2/n^2).
\]
For example, the nonedge product expansion follows by writing it as
$\exp\{(1-\alpha)\log(1-u)+\alpha\log(1-v)\}$;
the first and second derivatives of the nonedge factors are uniformly
bounded for $u,v\le1/2$. Expanding $-\log(1-x)$ and summing gives
\begin{equation}
 D^*_{i,ab}
 =n\max_{0\le \alpha\le1}\sum_c\pi_{c,n}
  \{(1-\alpha)P_{ac}+\alpha P_{bc}-P_{ac}^{1-\alpha}P_{bc}^\alpha\}
  +O((\log n)^2/n+\log n/n),\label{eq:chernoff-expansion}
\end{equation}
where $\pi_{c,n}=n_c/n$. The $O(\log n/n)$ term accounts for the
excluded node. Since $(n/\log n)P$ is fixed and $\pi_{c,n}\to\pi_c>0$,
the displayed expression divided by $\log n$ has a finite limit.
Distinct normalized block rows make this limit positive: taking
$\alpha=1/2$ gives one half the positive weighted sum of the squared
differences of their square roots. Thus $D^*_{i,ab}\asymp \log n$
uniformly over the finitely many pairs and excluded nodes, and
$I_n\asymp \log n$.

The equal-prior Bayes risk is
$R_{i,ab}^{\mathrm{or}}=\frac12\sum_x\min\{f_a(x),f_b(x)\}$.
Since $\min\{u,v\}\le u^{1-\alpha}v^\alpha$,
\begin{equation}
 R_{i,ab}^{\mathrm{or}}\le\tfrac12e^{-D^*_{i,ab}}.
 \label{eq:oracle-upper}
\end{equation}
For completeness, a logarithmic lower bound follows by exponential
tilting, without invoking any refined rate formula.  Let $\alpha_*$ maximize
$D_\alpha(p_{a,-i}\Vert p_{b,-i})$ and set
\[
 f_*(x)=e^{D^*_{i,ab}}f_a(x)^{1-\alpha_*}f_b(x)^{\alpha_*},
 \qquad L(x)=\log\frac{f_a(x)}{f_b(x)}.
\]
All masses are positive, the two distributions differ, and the endpoint
derivatives have opposite signs.  Thus $\alpha_*\in(0,1)$, and differentiating
the normalization gives $\mathbb E_*L=0$.  The tilted incident edges
remain independent, with probabilities
\[
 p_{*,j}=
 \frac{p_{aj}^{1-\alpha_*}p_{bj}^{\alpha_*}}
 {p_{aj}^{1-\alpha_*}p_{bj}^{\alpha_*}+(1-p_{aj})^{1-\alpha_*}(1-p_{bj})^{\alpha_*}}
 \le C\log n/n.
\]
The coefficient of each edge in $L$ is bounded by a constant depending
only on $b_-,b_+$; hence $\operatorname{Var}_*(L)\le C\log n$.
Taking $h_n=(\log n)^{2/3}$, Chebyshev's inequality gives
$\Pr_*(|L|\le h_n)=1-o(1)$.  The change of measure yields
\begin{align*}
 R_{i,ab}^{\mathrm{or}}
 &=\tfrac12e^{-D^*_{i,ab}}
   \mathbb E_*\min\{e^{\alpha_*L},e^{-(1-\alpha_*)L}\}\\
 &\ge\tfrac12e^{-D^*_{i,ab}-h_n}\Pr_*(|L|\le h_n).
\end{align*}
Combining this with \eqref{eq:oracle-upper} proves
$\log R_{i,ab}^{\mathrm{or}}=-D^*_{i,ab}+o(\log n)$, and hence
\eqref{eq:oracle-risk}.  For priors $\nu,1-\nu$ bounded away from zero,
$\min\{\nu f_a,(1-\nu)f_b\}$ is bounded above and below by constant
multiples of $\min\{f_a,f_b\}$, so its leading exponent is unchanged.
\end{proof}

\subsection{Transfer of the leading exponent}

\begin{proof}[Proof of Theorem~\ref{thm:exponent-transfer}]
Suppose the true label of node $i$ is $a$.  If $G_i$ occurs and the
classifier prefers $b\ne a$, then
$\widehat\ell_i(a)-\widehat\ell_i(b)\le0$, and consequently
$L_i^{ab}\le\tau_n$.  For every $\alpha\in[0,1]$,
\[
 \Pr_a\{L_i^{ab}\le\tau_n\}
 \le e^{\alpha\tau_n}\mathbb E_a e^{-\alpha L_i^{ab}}
 =\exp\{\alpha\tau_n-D_\alpha(p_{a,-i}\Vert p_{b,-i})\}.
\]
Selecting a Chernoff maximizer and using $\alpha\le1$ gives
\[
 \Pr\{\widehat z_i\ne z_i\}
 \le\delta_n+\sum_{b\ne a}\exp\{\tau_n-D^*_{i,ab}\}
 \le\delta_n+(K-1)\exp\{-I_n+o(\log n)\}.
\]
Since $I_n\asymp \log n$, we have $o(\log n)=o(I_n)$, and
\eqref{eq:transfer-tail} makes the first term no larger in leading
exponential order.  Averaging over nodes and using that the
permutation-invariant error is no larger than the error under any one
common alignment proves \eqref{eq:conditional-rate}.

The expected number of errors under that alignment is at most
$n\exp\{-(1+o(1))I_n\}$. If
$\liminf_{n\to\infty}I_n/\log n>1$, Markov's inequality gives vanishing
probability of even one error, which proves the exact recovery assertion.
No conclusion when $I_n/\log n\to1$ follows from this argument.
\end{proof}

The matching binary lower bound mentioned after the theorem has a direct
decision-theoretic justification.  Conditional on all other labels, the
rest of the graph is independent of the incident-edge row and has the
same distribution under $a$ and $b$.  Revealing the other labels and
true parameters cannot increase the Bayes risk.  A function of the
graph's spectral summary is also a function of the full graph; its risk
in that binary experiment is therefore at least
$R_{i,ab}^{\mathrm{or}}$.  This argument does not, by itself, assert a
global minimax lower bound over an unspecified parameter space.


\subsection{Full-graph reconstruction and local empirical EM}
\label{app:local-em}

\begin{proof}[Proof of Lemma~\ref{lem:full-spectral}]
The true no-self-loop mean $\E[A\mid z]$ leaves $\operatorname{col}(Z)$ invariant. In the orthonormal membership basis $ZN^{-1/2}$ its restriction is
$N^{1/2}PN^{1/2}-\operatorname{diag}(P_{11},\ldots,P_{KK})$.
Nonsingularity of the fixed matrix $(n/\log n)P$ and positive limiting proportions imply that its $K$ eigenvalues have absolute values between $\kappa_0\log n$ and $C \log n$ for some $\kappa_0>0$ and all large $n$. On the orthogonal complement, the eigenvalues are $-P_{aa}=O(\log n/n)$. Let $U_*$ and $\Lambda_*$ contain these $K$ signal eigenpairs; then
$U_*U_*^\top=\Pi_*$ and $\|U_*\|_{2\to\infty}=O(n^{-1/2})$.

Here $\|M\|_{2\to\infty}=\max_i\|M_i\|_2$. We spell out how the cited entrywise theorem applies. Fix any desired polynomial confidence power. Lei--Rinaldo's spectral norm theorem~\citep[Theorem 5.2]{leirinaldo2015}, with its confidence power chosen arbitrarily large, gives $\|A-\E[A\mid z]\|_{\mathrm{op}}\le C_H\sqrt{\log n}$ outside $O(n^{-H})$. Its upper degree parameter can be taken as $\max\{1,b_+\}\log n$. The signal separation is of order $\log n$, and the eigenvalue condition number is bounded. In the notation of Abbe et al.~\citep[Theorem 2.1]{abbe2020}, take $\gamma=C_H/\sqrt{\log n}$. Their incoherence assumption holds because $\|\E[A\mid z]\|_{2\to\infty}=O(\log n/\sqrt n)$; their row/column independence assumption follows from the independent upper-triangular edges.

For row concentration, their Bernoulli Lemma 7 applies with upper probability $b_+\log n/n$. Choosing its tail parameter large enough gives row failure $O(n^{-H-2})$. For matrices with a fixed number of columns, applying it to the columns and using monotonicity of $\varphi_H$ and of $\varphi_H(x)/x$ to replace the column norms by the Frobenius and maximum-row norms gives the required function
$\varphi_H(x)=C_H/(1\vee\log(1/x))$ for $x>0$, with $\varphi_H(0)=0$.
The constants absorb the fixed number of columns and the ratio of the signal gap to $\log n$. This function is continuous and nondecreasing, $\varphi_H(x)/x$ is nonincreasing, and both $\gamma$ and $\varphi_H(\gamma)$ tend to zero. Thus the remaining smallness condition in that theorem holds for all sufficiently large $n$.

Apply the theorem separately to the positive and negative signal groups, each a contiguous group of eigenvalues, and concatenate their orthogonal alignments. Repeated eigenvalues within either group cause no difficulty. The theorem allows a non-low-rank mean, so the excluded diagonal has not been silently discarded. The norm event also ensures that the empirical largest-absolute-value eigenpairs are precisely these signal groups. We obtain an orthogonal $R$, $V=UR$, and deterministic $\nu_n\to0$ such that
\[
 \|V\|_{2\to\infty}\le C_H/\sqrt n,\qquad
 \|V-AU_*\Lambda_*^{-1}\|_{2\to\infty}\le\nu_n/\sqrt n.
\]
The same polar alignments satisfy the Davis--Kahan bound
$\|V-U_*\|_{\mathrm{op}}\le C_H/\sqrt{\log n}$.
Put $M=R^\top\Lambda R=V^\top A V$. Since
$\|\E[A\mid z]\|_{\mathrm{op}}=O(\log n)$,
\[
 \|M-\Lambda_*\|_{\mathrm{op}}
 \le\|A-\E[A\mid z]\|_{\mathrm{op}}
   +2\|\E[A\mid z]\|_{\mathrm{op}}\|V-U_*\|_{\mathrm{op}}
 \le C_H\sqrt{\log n}.
\]
A degree union bound, with the degree-tail constant chosen for the same confidence power, gives $\max_i\deg(i)\le C_H\log n$.

Set $E=V-AU_*\Lambda_*^{-1}$. The exact row decomposition is
\[
 (\widetilde A-A\Pi_*)_i
 =E_iMV^\top+
 (AU_*)_i\{\Lambda_*^{-1}MV^\top-U_*^\top\}.
\]
The bracket has operator norm $O((\log n)^{-1/2})$, while
$\|(AU_*)_i\|_2\le\deg(i)\|U_*\|_{2\to\infty}\le C_H\log n/\sqrt n$.
Consequently
\[
 \|\widetilde A-A\Pi_*\|_{2\to\infty}
 \le C_H(\log n\nu_n+\sqrt{\log n})/\sqrt n
 =o(\log n/\sqrt n).
\]
Cauchy--Schwarz yields the required uniform row $\ell_1$ bound. Entrywise clipping is 1-Lipschitz, so the bound also holds for $q-q^0$. Moreover,
$|\widetilde A_{ij}|\le\|U_i\|_2\|\Lambda\|_{\mathrm{op}}\|U_j\|_2
\le C_H\log n/n$, proving the profile upper bound.

For the average approximation, define the block counts
$N_{ic}=\sum_{j:z_j=c}A_{ij}$. Because $c_0\log n/n<P_{ac}$, clipping toward $P_{ac}$ cannot increase the coordinate error:
\[
 \|q_i^0-p_{z_i,*}\|_1
 \le\sum_c|N_{ic}-n_cP_{z_i,c}|.
\]
Let $F$ be the sum on the right over all $i$. Its expectation is at most $Cn\sqrt{\log n}+C\log n$; the latter term accounts for the missing diagonal. Changing one independent edge changes $F$ by at most two. The bounded-differences inequality therefore gives
\[
 \Pr\{F-\E F>n(\log n)^{3/4}\}\le e^{-c (\log n)^{3/2}}.
\]
In logarithmic density this is smaller than any fixed inverse polynomial. Combining this average bound with the uniform $q-q^0$ bound, and enlarging a deterministic vanishing $a_n$, proves all assertions.
\end{proof}

\paragraph{Clipped oracle comparisons.}
For later use, fix a true class $a$ and a competitor $b$. Write
$v_c=\log(P_{ac}/P_{bc})$, $N'_{ic}=\max\{N_{ic},n_cc_0\log n/n\}$, and
\[
 L_{i,\mathrm{clip}}^{ab}
 =\sum_cN'_{ic}v_c-\sum_c n_c(P_{ac}-P_{bc}).
\]
This is a comparison score for the analysis, not a computable oracle supplied to the algorithm. For any $\alpha\in[0,1]$, a coordinate with $v_c\ge0$ has a clipped exponential moment no larger than its unclipped moment. If $v_c<0$, put $\beta=-\alpha v_c\ge0$ and $t_c=n_cc_0\log n/n$. For all sufficiently large $n$, $t_c\le\E_a N_{ic}$, including the excluded self-loop, because $c_0<\min_{a,c}nP_{ac}/\log n$. Jensen's inequality implies
\[
 \E_a e^{\beta\max\{N_{ic},t_c\}}
 \le\E_a e^{\beta N_{ic}}+e^{\beta t_c}
 \le2\E_a e^{\beta N_{ic}}.
\]
The counts over different column blocks are independent for a fixed node. Thus clipping costs at most a factor $2^K$ in the Chernoff moment. A direct binomial moment expansion, including an $O(\log n/n)$ diagonal correction, gives
\[
 \log\E_a e^{-\alpha L_{i,\mathrm{raw}}^{ab}}
 =-D_\alpha(p_{a,-i}\Vert p_{b,-i})+O((\log n)^2/n+\log n/n),
\]
where the raw score replaces $N'_{ic}$ by $N_{ic}$.
For any deterministic $\tau_n=o(\log n)$, exponential Markov inequality at the maximizing $\alpha$ yields
\begin{equation}
 \Pr_a\{L_{i,\mathrm{clip}}^{ab}\le\tau_n\}
 \le2^K\exp\{-D^*_{i,ab}+o(\log n)\}.
 \label{eq:clipped-chernoff}
\end{equation}
The argument controls the error probability directly. It does not assume that clipping changes every rare-row score by $o(\log n)$.

\begin{proof}[Proof of Theorem~\ref{thm:local-spectral-em}]
Choose the confidence power $H>2+\limsup_{n\to\infty}I_n/\log n$ in Lemma~\ref{lem:full-spectral} and work on its event. Constants below may depend on the fixed model and floor but not on the iteration. Suppose a responsibility matrix $R$ has aligned average error at most $\gamma$, with $\gamma$ below a sufficiently small fixed constant. Put
$N_a=\sum_iR_{ia}$ and $\bar q_a=n_a^{-1}\sum_{i:z_i=a}q_i$.
Then $|N_a-n_a|\le n\gamma$ and $N_a\ge\pi_{\min}n/2$, where $\pi_{\min}$ is a fixed lower bound on the community proportions for large $n$. The average-profile bound gives
$\|\bar q_a-p_{a*}\|_1\le C a_n\log n$.
The M-step identity gives, deterministically,
\[
 \|\widehat p_a-\bar q_a\|_1
 \le\frac{\sum_i|R_{ia}-Z_{ia}|\|q_i\|_1}{N_a}
       +\left|\frac{n_a}{N_a}-1\right|\|\bar q_a\|_1
 \le C\gamma \log n.
\]
Hence
\begin{equation}
 \max_a\|\widehat p_a-p_{a*}\|_1
 \le C(\gamma+a_n)\log n,\qquad
 c_0\log n/n\le\widehat p_{aj}\le C_H\log n/n.
 \label{eq:local-profile-fit}
\end{equation}
The logarithm is Lipschitz on this common positive range, so its vector $\ell_1$ error is at most $Cn/\log n$ times the profile error. Since $q_{ij}\le C_H\log n/n$, both the fitted linear contrasts and their offsets differ from the true-profile working contrasts by at most $C(\gamma+a_n)\log n$.
The fitted log-weight ratios are bounded: in the unconstrained update $N_a/n$ is bounded away from zero; for a fixed positive floor this follows directly from feasibility.

Combining the uniform row $q-q^0$ bound with these parameter bounds gives, for every node and pair,
\begin{equation}
 |\widehat\ell_i(a)-\widehat\ell_i(b)-L_{i,\mathrm{clip}}^{ab}|
 \le C(\gamma+a_n)\log n+C.
 \label{eq:local-score-fit}
\end{equation}
No independence of $R$, $q$, and the incident edges has been used. With $\gamma=\gamma_n\to0$, a wrong classification implies
$L_{i,\mathrm{clip}}^{ab}\le o(\log n)$ for some competitor, apart from the initialization or graph exceptional events. Equation~\eqref{eq:clipped-chernoff} and a finite union bound prove the one-update rate.

We also verify that further EM updates remain in the required region. Let
\[
 \kappa=\liminf_n\min_{a\ne b}
 (\log n)^{-1}D(p_{a*}\Vert p_{b*})>0.
\]
Choose a small fixed basin radius $\gamma_0>0$ so that the error bound from \eqref{eq:local-profile-fit} for $\gamma\le\gamma_0$ is less than $\kappa \log n/8$ in each score comparison for large $n$. Also choose a fixed small $t_0>0$ so that a row with
$\|q_i-p_{z_i,*}\|_1\le t_0\log n$ has an observation-contrast error less than $\kappa \log n/8$.
At the population observation $q_i=p_{z_i,*}$, the working margin equals $D(p_{z_i,*}\Vert p_{b*})\ge\kappa \log n/2$ for large $n$. The specified row deviation costs at most $\kappa \log n/8$, and the fitted-parameter and bounded weight corrections together cost at most $\kappa \log n/8$, leaving a correct-class margin of at least $\kappa \log n/4$.
The average-profile bound shows that the remaining rows constitute at most $a_n/t_0$ of all nodes. Thus the next E-step has aligned average responsibility error at most
\[
 b_n=2a_n/t_0+2(K-1)e^{-\kappa \log n/4}\longrightarrow0.
\]
This bound holds uniformly for every $R$ in the fixed basin. Eventually $b_n<\gamma_0$, so the basin is invariant. Starting with error $\gamma_n\to0$, every later responsibility matrix has error at most $\max\{\gamma_n,b_n\}=o(1)$.
Equation~\eqref{eq:local-score-fit} therefore holds with one deterministic $o(\log n)$ bound for every later fitted score on the same event. This includes a data-dependent finite stopping time and needs no union bound over iterations.

Finally add the initialization failure $\delta_n$ and the graph failure $O(n^{-H})$ to the per-node Chernoff bound and average over nodes. Since $H>2+\limsup_{n\to\infty}I_n/\log n$, the graph failure is negligible compared with $e^{-I_n}$, and the assumed initialization reliability preserves the leading exponent. If $\liminf_{n\to\infty}I_n/\log n>1$, multiplying the expected error by $n$ and applying Markov's inequality gives exact recovery with probability tending to one.
\end{proof}

\begin{proof}[Proof of Corollary~\ref{cor:provable-warmstart}]
The spectral norm and Davis--Kahan bounds above give an orthogonal alignment with
$\|U-U_*R^\top\|_F^2=O((\log n)^{-1})$ outside $O(n^{-H})$.
The $K$ distinct population row centers have squared separation at least $2/n$, since the rows of $ZN^{-1/2}$ from different groups are orthogonal, with squared norms $1/n_a$ and $1/n_b$.
A fixed-factor approximate $K$-means solution has total squared error at most a constant times $\|U-U_*R^\top\|_F^2$. The deterministic center-matching argument (also used in Lei--Rinaldo's clustering bound) then gives $\Mis(\widetilde z,z)=O((\log n)^{-1})$: the total squared discrepancy of fitted centers from the population rows is $O((\log n)^{-1})$; positive community sizes force one fitted center near each distinct population center; each incorrectly assigned row contributes a constant multiple of $1/n$ to that discrepancy. Hard responsibilities consequently have aligned average error $2\Mis(\widetilde z,z)=O((\log n)^{-1})$.

For a concrete randomized implementation, the $K$-means++ expected squared-error bound is $C_K=8(\log K+2)$ times the optimum~\citep{arthur2007}. Conditional on $U$, Markov's inequality bounds the chance of exceeding $2C_K$ times the optimum by $1/2$ for each independent seeding. Taking at least $(H+2)\log_2 n$ independent seedings and selecting the smallest squared-error potential makes that failure at most $n^{-H-2}$. Any subsequent Lloyd steps can only lower this potential. The zero-optimum case is recovered exactly by the seeding. For fixed $K$ this costs polynomial time and uses only $U$. Choose $H>\limsup_{n\to\infty}I_n/\log n$, combine this event with the spectral event, and apply Theorem~\ref{thm:local-spectral-em}.
\end{proof}


\subsection{End-to-end likelihood initialization}
\label{app:end-to-end}

\paragraph{Proof of Theorem~\ref{thm:two-community-growing}.}
Fix $H>2+\limsup_{n\to\infty}I_n/\log n$ and work on the event of Lemma~\ref{lem:full-spectral}. All constants below may depend on $H$ and the fixed model. Enlarge its deterministic $a_n\to0$ if necessary to bound both displayed profile errors. Put $\pi_{a,n}=n_a/n$,
\[
 m_b=\sum_{a=1}^2\pi_{a,n}P_{ab},\qquad
 D_{\pi_n}(x\Vert y)=
 \sum_b\pi_{b,n}\{x_b\log(x_b/y_b)-x_b+y_b\}.
\]
The block values of $q_c^0$ are
\[
 v_{cb}=\max\left\{\frac{\sum_{j:z_j=b}A_{cj}}{n_b},
                         \frac{c_0\log n}{n}\right\},
 \qquad q_{cj}^0=v_{c,z_j}.
\]
This is only an analysis variable. In particular $v_c$ may be a rare extreme candidate, and we do not assume it is near its own community profile. Let
\[
 h_a(v)=\frac{n}{\log n}\sum_b\pi_{b,n}
 \{P_{ab}\log(v_b/m_b)-(v_b-m_b)\},\qquad
 \Phi_n(v)=\sum_a\pi_{a,n}[h_a(v)]_+.
\]
The global profile mean $\overline q$ differs from $(m_{z_j})_{j=1}^n$ in $\ell_1$ norm by at most $a_n\log n$. The uniform bound on $q_c-q_c^0$ and the common positive box give
\[
 \left|[\ell_i(q_c)-\ell_i(\overline q)]-\log nh_{z_i}(v_c)\right|
 \le C\{\|q_i-p_{z_i,*}\|_1+a_n\log n\}.
\]
Here the parameter perturbation costs
$C\|q_c-q_c^0\|_1+C\|\overline q-(m_{z_j})_{j=1}^n\|_1$,
because $q_{ij}\le C \log n/n$ and $\log$ is $n/(c_0\log n)$-Lipschitz. The remaining observation perturbation multiplies a uniformly bounded log ratio. Since positive part is 1-Lipschitz, summing yields the simultaneous candidate bound
\[
 \sup_c\left|\frac{\Delta_2(c)}{n\log n}-\Phi_n(v_c)\right|\le Ca_n.
\]
Every true class has a candidate with $\|q_c-p_{a*}\|_1=o(\log n)$ by the average error bound and its positive proportion. The same is true of $q_c^0$, so
$(n/\log n)\sum_b\pi_{b,n}|v_{cb}-P_{ab}|=o(1)$.
The fixed normalized matrix and positive limiting proportions imply that
\[
 g_0=\lim_{n\to\infty}\max_a
 \frac{n\pi_{a,n}}{\log n}
 D_{\pi_n}\big((P_{ab})_{b=1}^2\Vert m\big)>0.
\]
Positivity follows because the two distinct normalized block rows cannot both equal their mean. Continuity on the normalized common positive box and the candidate bound imply that the selected candidate satisfies $\Phi_n(v_c)\ge g_0/2$ for large $n$.

The decisive two-class identity is
\[
 \sum_a\pi_{a,n}h_a(v)=-\frac{n}{\log n}D_{\pi_n}(m\Vert v)\le0.
\]
Together with $\Phi_n(v_c)\ge g_0/2$, it forces exactly one of $h_1(v_c),h_2(v_c)$ to be positive. Moreover
\[
 h_{\rm pos}(v_c)\ge g_0/2,\qquad
 h_{\rm neg}(v_c)\le-g_0/2,
\]
since each weighted positive or negative magnitude is at least $g_0/2$ and each proportion is at most one. The initial two mixture weights are equal, so they do not alter this score comparison.

Take $\tau_n=\sqrt{a_n}+(\log n)^{-1}$. At most $(a_n/\tau_n)n=o(n)$ observations have $\|q_i-p_{z_i,*}\|_1>\tau_n\log n$. Every other observation has the correct initial component margin at least $g_0\log n/4$ after the above score perturbations. After a common label alignment, the first E-step therefore satisfies
\[
 \frac1n\sum_{i,a}|R_{ia}-Z_{ia}|
 \le \frac{2a_n}{\tau_n}+2e^{-g_0\log n/4}=o(1).
\]
The graph event has failure $O(n^{-H})$, so these responsibilities meet the reliability condition in Theorem~\ref{thm:local-spectral-em}. The evaluated fitting loop always performs its first M-step before it may stop. That theorem gives the claimed exponent for the first resulting decode and all subsequent finite updates, including the evaluated stopping rule. Its exact-recovery implication also applies. \qed

\paragraph{Proof of Theorem~\ref{thm:certified-likelihood}.}
Again choose a fixed $H>2+\limsup_{n\to\infty}I_n/\log n$ and condition on the event in Lemma~\ref{lem:full-spectral}. Let $a_n\to0$ bound its average profile error, and write $c_0\log n/n\le q_{ij}\le C_H\log n/n$. True profiles have the same positive box after enlarging $C_H$. Full rank implies distinct block rows; their limiting positive proportions give constants $\kappa,s>0$ such that
\[
 \min_{a\ne b}D(p_{a*}\Vert p_{b*})\ge\kappa \log n,\qquad
 \min_{a\ne b}\|p_{a*}-p_{b*}\|_1\ge s\log n
\]
for all sufficiently large $n$.

\emph{Reliable seed coverage.}
On this box the two derivatives of the scalar deviance are $\log(x/y)$ and $1-x/y$, which are uniformly bounded. Consequently
\[
 \begin{split}
 0\le D(x\Vert y)&\le C\log n,\\
 |D(x\Vert y)-D(x'\Vert y')|
 &\le C(\|x-x'\|_1+\|y-y'\|_1).
 \end{split}
\]
Set $\tau_n=\sqrt{a_n}+(\log n)^{-1}$. Call row $i$ good when $\|q_i-p_{z_i,*}\|_1\le\tau_n\log n$; the remaining rows number at most $(a_n/\tau_n)n$. Conditional on previously selected seeds being good and representing distinct classes, every good row in a represented class has residual at most $2C\tau_n\log n$. Every good row in an unrepresented class has residual at least $\kappa \log n/2$, by the Lipschitz bound. Before $K$ seeds have been selected, an unrepresented class supplies at least $\pi_{\min}n/2$ good rows. Hence total sampling mass is at least $cn\log n$, whereas the mass of represented good rows and all bad rows is at most
\[
 C(\tau_n+a_n/\tau_n)n\log n.
\]
The first uniform seed is bad with probability at most $a_n/\tau_n$. A union bound over the fixed $K$ choices gives single-run failure at most
\[
 r_{H,n}\le C_K(\tau_n+a_n/\tau_n)=o(1).
\]
The all-zero fallback cannot occur on this successful path while a class remains unrepresented. Conditional on $q$, the restart streams are independent. Their joint failure probability is at most
\[
 r_{H,n}^{\lceil\log n\rceil}=n^{-\omega(1)}.
\]
This amplification requires no knowledge of $H$ or the separation constants.

\emph{Likelihood selection certifies the winning fit.}
Define the observation-only saturation constant and soft deviance
\[
 S_q=\sum_{i,j}(q_{ij}\log q_{ij}-q_{ij}),\qquad
 F(\widehat\pi,\widehat p)=
 -\sum_i\log\sum_k\widehat\pi_k
       e^{-D(q_i\Vert\widehat p_{k*})}.
\]
There is the exact identity
\[
 \mathcal L=S_q-F,\qquad
 F\ge\sum_i\min_kD(q_i\Vert\widehat p_{k*})\ge0.
\]
Observation Gamma terms, if restored, are common to every fit. A successful run starts with one good seed for each class and uniform weights. Thus
\[
 F_{\rm initial}
 \le C(a_n+\tau_n)n\log n+n\log K
 =:\epsilon_nn\log n,\qquad \epsilon_n\to0.
\]
Every exact EM update increases $\mathcal L$ and decreases $F$. Selecting the largest final $\mathcal L$ across all runs, including the growing endpoint, therefore gives
\[
 F_{\rm selected}\le\epsilon_nn\log n.
\]
This argument does not require identifying the successful run, fitting any run to stationarity, or proving anything about the growing branch's basin dynamics.

All candidate means remain in the same box, since they start at rows or the global mean of $q$ and each M-step forms a convex combination. Keeping the previous mean in a zero-mass component preserves the box too. Strong convexity of $x\log x-x$ on the box, followed by Cauchy--Schwarz, gives
\[
 D(x\Vert y)
 \ge\frac{n\|x-y\|_2^2}{2C_H\log n}
 \ge\frac{\|x-y\|_1^2}{2C_H\log n}.
\]
Consequently
\[
 \sum_i\min_k\|q_i-\widehat p_{k*}\|_1^2
 \le 2C_H\epsilon_nn(\log n)^2.
\]
Put $\eta_n=\sqrt{a_n}+\epsilon_n^{1/4}+(\log n)^{-1}$. The fractions of rows farther than $\eta_n\log n$ from their true profile, or from every fitted profile, are at most $a_n/\eta_n$ and $2C_H\epsilon_n/\eta_n^2$, respectively; both tend to zero. Each positive-proportion true class has a row avoiding both bad sets. Its fitted center is within $2\eta_n\log n$ of its true center. The separation $s\log n$ forces different classes to use different fitted centers; with exactly $K$ centers this gives a single label permutation satisfying
\[
 \max_a\|\widehat p_{a*}-p_{a*}\|_1\le2\eta_n\log n.
\]

\emph{The selected responsibilities are weakly consistent.}
For a row with $\|q_i-p_{z_i,*}\|_1\le\eta_n\log n$, every wrongly matched center is at $\ell_1$ distance at least $s\log n/2$ for large $n$, hence its deviance is at least
$c_*\log n$, with $c_*=s^2/(8C_H)>0$.
For the selected fit's exact E-step, write
\[
 R_{ik}=
 \frac{\widehat\pi_k e^{-D(q_i\Vert\widehat p_{k*})}}
 {\sum_l\widehat\pi_l e^{-D(q_i\Vert\widehat p_{l*})}}.
\]
The variational identity is
\[
 -\log\sum_k\widehat\pi_k e^{-D_{ik}}
 =\sum_kR_{ik}D_{ik}
   +\operatorname{KL}(R_i\Vert\widehat\pi)
 \ge\sum_kR_{ik}D_{ik}.
\]
Summing over the typical rows and bounding each remaining row's responsibility error by two gives
\[
 \frac1n\sum_{i,k}|R_{ik}-Z_{ik}|
 \le\frac{2a_n}{\eta_n}+\frac{2\epsilon_n}{c_*}
 =:\gamma_n\longrightarrow0
\]
after the above permutation. The identity remains valid with zero weights using the usual zero-mass conventions, and also with a feasible fixed weight floor.

This weak-start conclusion fails only if the graph event fails or every restart fails, a probability at most $O(n^{-H})+n^{-\omega(1)}$. Apply Theorem~\ref{thm:local-spectral-em} to the mandatory final M-step. Since $H>2+\limsup_{n\to\infty}I_n/\log n$, the initialization failures are negligible at the target exponent scale. This proves the complete decoder's expected-error bound and its exact-recovery consequence. Prescribed polynomial EM iteration limits and $\lceil\log n\rceil$ restarts give polynomial computational cost. \qed

\subsection{Row-independent projection errors}

\begin{proof}[Proof of Theorem~\ref{thm:independent-projection}]
Condition on the data generating $\Pi_i$ and on $H_i$.  Work in the
$n-1$ coordinates other than $i$, and write $\Pi_*=\Pi_{*,i}$,
$\Pi=\Pi_i$, $w=w_{-i}^{ab}$, $p_j=\mathbb E A_{ij}$, and
$e=(\Pi-\mathrm{Id})w$.  The fixed block assumptions imply
\[
 w\in\operatorname{col}(\Pi_*),\quad
 p\in\operatorname{col}(\Pi_*),\quad
 \|w\|_2\le C\sqrt n,\quad \|w\|_\infty\le C,
 \quad \|p\|_2\le C\log n/\sqrt n.
\]
The constants are uniform over the finitely many contrasts and true
labels.  For orthogonal projectors,
\[
 \|\Pi_*(\mathrm{Id}-\Pi)\Pi_*\|_{\mathrm{op}}
 =\|(\mathrm{Id}-\Pi)\Pi_*\|_{\mathrm{op}}^2\le\eta_n^2,
\]
because the left matrix is
$[(\mathrm{Id}-\Pi)\Pi_*]^\top[(\mathrm{Id}-\Pi)\Pi_*]$.  It follows that the mean projection
error has the quadratic bound
\begin{equation}
 |p^\top e|
 \le\|p\|_2\|w\|_2\eta_n^2
 \le C\log n\eta_n^2.\label{eq:quadratic-bias}
\end{equation}
Also,
$\|e\|_2\le C\sqrt n\eta_n$.  By Cauchy--Schwarz and projector
leverage,
\[
 |(\Pi w)_j|\le\|\Pi\mathbf e_j\|_2\|w\|_2
 =\sqrt{\Pi_{jj}}\|w\|_2\le C\sqrt L,
\]
so $\|e\|_\infty\le M$ for a fixed positive constant $M$.

Conditional independence makes
$X_j=(A_{ij}-p_j)e_j$ independent and centered.  They satisfy
$|X_j|\le M$ and
\begin{equation}
 V=\sum_j\mathbb E X_j^2
 \le b_+\log n/n\|e\|_2^2\le C\log n\eta_n^2.
 \label{eq:projection-variance}
\end{equation}
Here a Bennett bound is stronger than the simpler Bernstein bound at
the required scale.  We include its short derivation.  For a centered
$X$ with $|X|\le M$ and $\lambda\ge0$, Taylor expansion and
$|X|^k\le M^{k-2}X^2$ imply
\[
 \mathbb E e^{\lambda X}
 \le1+\frac{\mathbb E X^2}{M^2}
                (e^{\lambda M}-1-\lambda M).
\]
After multiplying moment-generating functions and using $\log(1+x)\le x$,
exponential Markov inequality optimized in $\lambda$ gives
\begin{equation}
 \Pr\{|\textstyle\sum_jX_j|>t\mid \Pi\}
 \le2\exp\left\{-\frac{V}{M^2}
 h\left(\frac{Mt}{V}\right)\right\},\qquad
 h(u)=(1+u)\log(1+u)-u,\label{eq:bennett}
\end{equation}
with the evident interpretation when $V=0$.

To choose a uniform vanishing error scale, put
$a_n=\max\{\eta_n^2,e^{-\log n}\}$,
$s_n=\log(1/a_n)$.
Both $a_n\to0$ and $s_n\to\infty$.  Take $t=\log n/\sqrt{s_n}$.
The bound \eqref{eq:bennett} remains valid with $V$ replaced by the
upper bound $C\log na_n$, since its exponent is nonincreasing in the
variance argument.  As $u\to\infty$, $h(u)\sim u\log u$, so
\[
 \frac{C\log na_n}{M^2}
 h\left(\frac{M}{Ca_n\sqrt{s_n}}\right)
 =\frac{\log n}{M\sqrt{s_n}}
   \{s_n-\tfrac12\log s_n+O(1)\}(1+o(1))
 \ge c \log n\sqrt{s_n}
\]
for all sufficiently large $n$.  Combining this with
\eqref{eq:quadratic-bias} gives, for
$\tau_n=\log n/\sqrt{s_n}+C\log n\eta_n^2=o(\log n)$,
\[
 \Pr\{|A_{i,-i}(\Pi-\mathrm{Id})w|>\tau_n\mid \Pi,H_i\}
 \le2e^{-c \log n\sqrt{s_n}}.
\]
Union bound over the finitely many contrasts and addition of
$\Pr(H_i^c)$ prove \eqref{eq:independent-projection-tail} with
$h_n=c\sqrt{s_n}\to\infty$.  The last claim is an application of
Theorem~\ref{thm:exponent-transfer}.
\end{proof}

\section{Proofs of initialization and optimization properties}
\label{app:optimization-proofs}

\subsection{Frozen-candidate greedy gains}

\begin{proposition}[Greedy gain guarantee]
\label{prop:greedy-gain}
Let $s_{ic}$ be fixed finite candidate scores and let $S_0$ be a fixed
nonempty seed set.  For candidate sets $T$, define
\[
 G(T)=\sum_i\max_{c\in S_0\cup T}s_{ic}
      -\sum_i\max_{c\in S_0}s_{ic}.
\]
Then $G$ is normalized, monotone, and submodular.  If $m\ge1$ greedy
steps select a set $T_m$ and $T^*$ maximizes $G$ over sets of at most
$m$ candidates, then
\[
 G(T_m)\ge\{1-(1-1/m)^m\}G(T^*)\ge(1-e^{-1})G(T^*).
\]
\end{proposition}

\begin{proof}
Normalization and monotonicity follow immediately from the definition.
The marginal gain of candidate $c$ is
\[
 G(T\cup\{c\})-G(T)
 =\sum_i\left(s_{ic}-\max_{b\in S_0\cup T}s_{ib}\right)_+.
\]
If $T\subseteq T'$, each maximum in the right side increases, so the
marginal gain decreases.  This proves submodularity.
For the greedy set $T_t$ at time $t$, monotonicity and telescoping
diminishing returns give
\begin{align*}
 G(T^*)-G(T_t)
 &\le G(T_t\cup T^*)-G(T_t)\\
 &\le\sum_{c\in T^*\setminus T_t}
        \{G(T_t\cup\{c\})-G(T_t)\}\\
 &\le m\{G(T_{t+1})-G(T_t)\}.
\end{align*}
Starting from $G(T_0)=0$ and iterating this inequality proves the
claimed factor.  The final inequality follows from
$(1-1/m)^m\le e^{-1}$, including the case $m=1$.
\end{proof}

This is the standard cardinality-constrained submodular greedy
guarantee~\cite{nemhauser1978}.  It controls added explanation relative
to the fixed initial seed, not the multiplicative approximation of a
possibly negative total log score.  It does not apply unchanged when
intermediate EM steps alter the candidate scores.

\subsection{Noiseless profile selection}

\begin{proposition}[Population profile selection]
\label{prop:noiseless-seeding}
Suppose $q_i=\theta_{z_i}$, where the $K$ positive profile vectors
$\theta_1,\ldots,\theta_K$ are pairwise distinct and every class is
nonempty.  The candidate pool consists of these node profiles.  From
any first seed belonging to this candidate pool, $K-1$ subsequent global-gain greedy steps select all
$K$ different profiles.  Maximum working-score classification using
those profiles recovers $z$ up to permutation.
\end{proposition}

\begin{proof}
For $x,\widehat p>0$ coordinatewise, write
$\ell_x(\widehat p)=\sum_jx_j\log\widehat p_j-\sum_j\widehat p_j$.  Then
\[
 \ell_x(x)-\ell_x(\widehat p)
 =\sum_j\{x_j\log(x_j/\widehat p_j)-x_j+\widehat p_j\}
 =D(x\Vert\widehat p)\ge0,
\]
with equality precisely when $x=\widehat p$; the inequality follows from
$\log u\le u-1$.  If a profile has already been selected, selecting a
duplicate has zero gain.  If profile $\theta_b$ has not been selected,
then each node in class $b$ gains
$\min_{a\in S}D(\theta_b\Vert\theta_a)>0$
when it is added.  Therefore some unselected profile has strictly
positive gain, and greedy cannot select a duplicate before all distinct
profiles have been selected.  The same strict divergence inequality
makes each node's own profile its unique best score after selection.
\end{proof}

The proposition concerns the frozen node-profile selection control in an abstract noiseless profile model. It does not directly cover Algorithm~1, whose initial mean need not be a node profile and whose intermediate EM fits change the parameters.  It does
not replace a stability analysis for random spectral profiles, and it
does not erase the diagonal distinction in a no-self-loop SBM.

\subsection{Exact working EM and accepted replacements}

\begin{proposition}[Objective convergence of exact working EM]
\label{prop:em-monotonicity}
For fixed strictly positive profiles $q_1,\ldots,q_n$, let
\[
 \mathcal L(\widehat\pi,\widehat p)=\sum_i\log\sum_{a=1}^K
       \widehat\pi_a\exp\{\ell_i(\widehat p_a)\},\qquad
 \ell_i(\widehat p_a)=\sum_jq_{ij}\log\widehat p_{aj}-\sum_j\widehat p_{aj}.
\]
Starting with positive parameters and positive mixture weights, the
exact E-step and unconstrained exact M-step
\begin{align*}
 r_{ia}&=\frac{\widehat\pi_a e^{\ell_i(\widehat p_a)}}
                   {\sum_b\widehat\pi_b e^{\ell_i(\widehat p_b)}},\qquad
 N_a=\sum_i r_{ia},\\
 \widehat\pi_a^{\mathrm{new}}&=N_a/n,\qquad
 \widehat p_{aj}^{\mathrm{new}}=\frac{\sum_i r_{ia}q_{ij}}{N_a}
\end{align*}
make $\mathcal L$ nondecreasing.  Its numerical values converge to a
finite limit.  The same monotonicity statement holds for an exact
M-step maximizing over any stated constraint set containing the current
parameters, and for replacement proposals accepted only when their
final $\mathcal L$ improves.
\end{proposition}

\begin{proof}
For row-wise probability vectors $r_i$, define the variational objective
\[
 \mathcal F(r,\widehat\pi,\widehat p)=
 \sum_{i,a}r_{ia}\{\log\widehat\pi_a+\ell_i(\widehat p_a)-\log r_{ia}\}.
\]
The log-sum inequality gives
$\mathcal F(r,\widehat\pi,\widehat p)\le\mathcal L(\widehat\pi,\widehat p)$, with equality
for the E-step responsibilities.  For fixed $r$, maximizing
$\sum_aN_a\log\widehat\pi_a$ over the simplex gives $\widehat\pi_a=N_a/n$.
For each profile coordinate, maximizing
$(\sum_i r_{ia}q_{ij})\log\widehat p_{aj}-N_a\widehat p_{aj}$ gives the displayed
weighted mean.  These are interior maxima because all inputs and
responsibilities are positive.  Thus the E-step makes the variational
bound tight and the M-step raises it, proving monotonicity.

The saturated-score inequality used in the preceding proof implies
$\ell_i(\widehat p_a)\le\ell_i(q_i)$ for every $i,a$.  Since
$\sum_a\widehat\pi_a=1$,
\[
 \log\sum_a\widehat\pi_a e^{\ell_i(\widehat p_a)}\le\ell_i(q_i).
\]
Consequently $\mathcal L\le\sum_i\ell_i(q_i)<\infty$.
A nondecreasing real sequence bounded above has a finite limit.
For a constrained exact M-step the current parameters remain feasible,
so the same variational argument applies.  The replacement assertion
follows directly from its acceptance rule.
\end{proof}

This proposition proves convergence of objective values.  It proves
neither convergence to the global maximum nor recovery of the true
labels. The implementation uses an exact active-set maximizer on the
floored simplex, so its weight constraint is covered by the constrained
version of this proposition. Its finite stopping tolerance and
floating-point accuracy are assessed separately in Section~\ref{sec:experiments}.

\begin{proposition}[Profile hull and spectral dimension]
\label{prop:profile-hull}
Every responsibility-weighted profile M-step produces a profile in the convex hull of the rows of $q$. If the unmodified reconstruction $q=U\Lambda U^\top$ is used and the resulting means are positive, all fitted profiles lie in its row space, of dimension at most $K$. Entrywise positive clipping can increase this dimension.
\end{proposition}
\begin{proof}
The coefficients $R_{ia}/N_a$ are nonnegative and sum to one, so the update is a row convex combination. Every such combination also lies in the row space. The unmodified reconstruction has rank at most $K$. The clipping operation is nonlinear and has no corresponding rank-preservation identity.
\end{proof}
